\documentclass[11pt]{article}
\usepackage[T1]{fontenc}
\usepackage[utf8]{inputenc}
\usepackage{lmodern,amsmath,amssymb,amsthm}
\usepackage[margin=1in]{geometry}
\usepackage[hidelinks]{hyperref}
\newtheorem{lemma}{Lemma}\newtheorem{theorem}[lemma]{Theorem}
\newcommand{\R}{\mathbb R}\newcommand{\Q}{\mathbb Q}
\newcommand{\Z}{\mathbb Z}\newcommand{\N}{\mathbb N}
\newcommand{\C}{\mathbb C}\newcommand{\DR}{\mathrm{DR}}
\title{Degree conversion and the finite rank iteration}
\author{Author proposal: /root/astra\_v14\_baseline\_theta}
\date{7 November 2026}
\begin{document}\maketitle

\noindent\textbf{Status.} This is a bounded author proof awaiting independent
review. It proves the degree conversions and the assembly of the rank
steps from the precisely stated one-step geometric/frequency interface.
It does not prove the analytic sunflower or commutator hypotheses of
GTZ Sections 10--11. The complete inverse theorem remains open.
Every internal nilcharacter below has exact unit norm and vertical
equivariance for all limit elements under the natural ultrapower
interpretation. The main manuscript is not modified.

\section{The actual degree-to-rank conversion}

Let $G/\Gamma$ be a connected simply connected nilpotent Lie group with
lattice and rational degree filtration $(G_i)_{i\geq0}$ supported on
$0\leq i\leq m$. Apply accepted bounded tautification first, so
$G_0=G_1=G$ and polynomial orbits may be normalized at zero without
changing the represented sequence. Write $\mathfrak g_i=\log G_i$.
All the vector spaces here are rational for the rational structure
determined by $\Gamma$.

For $(d,r)\in\DR=\{(d,r)\in\N^2:0\leq r\leq d\}$ with $d>0$ define
$\mathfrak h_{(d,r)}$ to be the real span of all iterated Lie brackets
of length $\ell\geq1$ with entries $X_j\in\mathfrak g_{i_j}$,
$i_j\geq1$, for which
\begin{equation}\label{eq:DR}
 w:=\sum_{j=1}^{\ell}i_j>d,
 \qquad\hbox{or}\qquad
 w=d\ \hbox{and}\ \ell\geq\max(r,1).
\end{equation}
Take every bracketing and set
$\mathfrak h_{(0,0)}=\mathfrak g$, $H_{(d,r)}=\exp\mathfrak h_{(d,r)}$.
Only finitely many nonzero brackets occur, since their weight exceeds
$m$ when they can no longer contribute and every entry has positive
weight.

\begin{lemma}\label{lem:DR}
These are connected rational groups forming a taut degree--rank
filtration supported on $(m,m)$. For each $d>0$,
$H_{(d,0)}=H_{(d,1)}=G_d$. Hence the original degree-polynomial orbits
are exactly the degree--rank polynomial orbits for this refinement.
A degree-$m$ nilcharacter is a degree--rank $(m,m)$ nilcharacter
with the same function and with the top character restricted.
\end{lemma}
\begin{proof}
All bracket spans in \eqref{eq:DR} are rational, because brackets of
rational vectors are rational. Bracketing with $\mathfrak g_1$ raises
weight by at least one, so each span is an ideal of $\mathfrak g$.
It is therefore a Lie subalgebra, and its exponential is a connected
normal subgroup, which is closed by the global exponential homeomorphism.
Choose a rational Mal'cev basis for $G/\Gamma$ compatible with its rational
structure. Each bracket span has a rational basis in these coordinates.
Its lower central series is rational, since brackets of rational vectors
are rational. Choose one standard $Q$ bounding the finitely many rational
structure constants and basis coefficients. Green--Tao, \emph{Annals of
Mathematics} 175 (2012), Definition 2.5 and Proposition A.10, applied to
this subgroup with its lower central filtration, supplies a Mal'cev basis
relative to its actual intersection with $\Gamma$. That intersection is
therefore discrete and cocompact. This application uses Lie-algebra
rationality in Definition 2.5; it does not assume the lattice conclusion.

The lexicographic nesting is immediate at a fixed $d$ as $r$ increases.
If $d$ increases, every new allowed weight was strictly above the old
$d$, so its span is contained in the old one as well. For two positive
indices $(d,r),(e,t)$, brackets either have weight greater than $d+e$,
or weight exactly $d+e$ and length at least
$\max(r,1)+\max(t,1)\geq\max(r+t,1)$. Thus the Lie bracket is in
$\mathfrak h_{(d+e,r+t)}$. The BCH commutator and its higher terms
lie there too, using the ideal property. At index zero use normality.
This verifies every pair of indices, not only equal indices.

For $r=0$ or $1$, every allowed bracket has weight at least $d$ and
thus lies in $\mathfrak g_d$. Conversely the length-one brackets with
$i_1=d$ include all of $\mathfrak g_d$. This proves
$H_{(d,0)}=H_{(d,1)}=G_d$. In particular $H_{(1,0)}=G$.
If $d>m$ all brackets vanish. At degree $m$ every nonzero bracket has
length at most its weight $m$, which gives the asserted support.

The domain filtration on $\Z$ has $\Z$ at indices $(0,0),(1,0)$
and zero at all other indices, including $(1,1)$, as in Definition 6.17.
The derivative criterion for an adapted
map therefore asks for its $a$th differences to lie in $H_{(a,0)}$;
these groups equal $G_a$. This proves the equality of the polynomial
classes directly, matching published (6.8). The new top
$H_{(m,m)}$ is a rational subgroup of $G_m$. Restricting the original
vertical character is integral on its lattice and preserves all exact
vertical identities and unit norm. No enlargement of its character
domain is asserted.
\end{proof}

\section{The total-degree conversion used in related applications}

For a rational multidegree filtration $(G_a)_{a\in\N^q}$ supported
on a finite downset $J$, first use bounded tautification to restrict
to $H=\langle G_{e_1},\ldots,G_{e_q}\rangle$. Let
\[
 K_0=H,\qquad
 \mathfrak k_i=\sum_{a\in J:\ |a|\geq i}\mathfrak g_a,
 \qquad K_i=\exp\mathfrak k_i\quad(i\geq1).
\]
Each $\mathfrak g_a$ is an ideal in $\mathfrak g_0$, so the sum is
a rational ideal. It equals the Lie algebra of the connected generated
group. The brackets satisfy
$[\mathfrak k_i,\mathfrak k_j]\subseteq\mathfrak k_{i+j}$,
because $[\mathfrak g_a,\mathfrak g_b]\subseteq\mathfrak g_{a+b}$
and $|a+b|=|a|+|b|$. The index-zero condition follows from normality.
Also $K_1=H$ and $K_i$ is trivial above $D=\max_{a\in J}|a|$.
The rational subgroup lemma verifies every lattice intersection.

For the diagonal restriction $n\mapsto g(n,\ldots,n)$ of a
multivariable polynomial orbit, use its Taylor expansion. A factor
with coefficient $u_a\in G_a$ has scalar exponent
$\prod_j\binom n{a_j}$ of ordinary degree $|a|$. Such a factor
is adapted to $K_\bullet$ since $u_a\in K_{|a|}$; its successive
differences have the same one-parameter coefficient and the required
scalar polynomial degree. The product is polynomial by B.4. This
checks the actual diagonal map rather than treating the domain
filtrations as identical. For a rectangle $J=\{a:a\leq d\}$,
the only index of total weight $|d|$ is $d$, so $K_{|d|}=G_d$.
A multidegree-$d$ nilcharacter therefore restricts diagonally to a
degree-$|d|$ nilcharacter with exactly the same top character.
For a nonrectangular downset only the nilsequence degree assertion
has been proved, and no unwarranted common top character is claimed.

\section{The precise one-step interface}

Fix $s\geq3$, write $m=s-1$, and let $1\leq R\leq m$.
Suppose $f$ is bounded on $[N]$, $H$ is a dense internal set of
shifts, $A(h,n)$ is a single unit multidegree $(1,m)$ nilcharacter,
and $(\chi_h)$ is an internal uniformly bounded family of unit
degree--rank $(m,R)$ nilcharacters. Suppose that for every retained
$h$, $\Delta_h f$ has positive degree-at-most-$m-1$ correlation with
$A(h,\cdot)\otimes\chi_h$.

The one-step premise used here is the following exact output, after
restricting to a dense internal subset and fixing uniform standard
bounds. It must be supplied by the corrected additive-quadruple result,
the actual Section 10 frequency representation, the two conclusions of
Theorem 11.1, accepted LP04 and the final Section 12 join:
\begin{enumerate}
\item There is a unit multidegree $(1,m)$ nilcharacter $Z(h,n)$,
defined globally, and a family of unit-vector sequences $\zeta_h(n)$
such that $Z(h,n)=\iota\zeta_h(n)$ for retained $h$, with one fixed
coordinate isometry $\iota$. Zero padding is allowed.
\item The entire matrix-valued function
$V_h=\chi_h\otimes\overline{\zeta_h}$ is a nilsequence of
degree--rank at most $(m,R-1)$, with standard uniform geometry and
Lipschitz bounds and internal parameter dependence.
\end{enumerate}
The second clause states a rank drop; for $R>1$ it does not state
an ordinary-degree drop. The final Section 12 constructor must choose
its test equivariant on the full ordinary top before invoking LP04's
comparison. Equivariance only on $G_{(m,R)}$ does not by itself
justify its larger multidegree top. Those obligations remain in the
separately authored geometric join and are not hidden in this premise.

\section{Recovering the next family without dividing by components}

The unit-vector identity gives, for every component $a$ and all relevant
integers,
\begin{equation}\label{eq:recovery}
 (\chi_h)_a(n)=\sum_b (\zeta_h)_b(n)
       \bigl((\chi_h)_a(n)\overline{(\zeta_h)_b(n)}\bigr).
\end{equation}
The terms in parentheses are components of $V_h$. Replace the first
factor by its corresponding coordinate in $Z(h,n)$ and form the unit
nilcharacter $A'(h,n)=A(h,n)\otimes Z(h,n)$ on the product model.
The product filtration has the same multidegree bound $(1,m)$.

Choose a scalar component and a degree-at-most-$m-1$ test witnessing
the incoming correlation, with fixed bounds and a fixed positive
threshold after the usual finite selection. Inserting
\eqref{eq:recovery} expresses it as a fixed finite sum of correlations
of $\Delta_h f$ against components of $A'(h,\cdot)$, the same
ordinary lower-degree test, and a component of $V_h$. At least one
summand has a positive fixed correlation. Finite pigeonhole restricts
to a dense subset with fixed component choices if necessary. No scalar
component is inverted, or assumed nonzero or of modulus one.

On each fixed rational degree--rank model for $V_h$, apply the
unit-character approximation construction to its top group at
$(m,R-1)$. That group is central after tautification, since bracketing
with $G_{(1,0)}$ raises the first coordinate above $m$. Its lattice
is rational. Smooth local sections and a finite unit partition on the
central torus produce exact-character unit-vector tests. The complex linear span of their components
is a unital conjugation-stable algebra separating points; Stone--Weierstrass
and a finite uniform net of the bounded Lipschitz test class give a
finite approximation by components of such tests. This is the same
explicit compact construction in the assembly's selection section,
now on the stated degree--rank model. It uses centrality and rationality,
not an identification of this top with the larger ordinary-degree top.

Choose the approximation error so that its correlation with the other
fixed bounded factors is less than half the surviving threshold.
Another finite selection supplies a unit degree--rank $(m,R-1)$
nilcharacter $\chi'_h$ for which $\Delta_h f$ has positive
degree-at-most-$m-1$ correlation with
$A'(h,\cdot)\otimes\chi'_h$. Finite choices and orbit choices are
made on representatives and then passed to the ultralimit, retaining
internal dependence and uniform bounds. Every test has its exact
standard vertical identity before extension, so the all-limit
convention is preserved. This is the output required by the next
rank step.

\section{Initialization, termination and the actual endpoint}

Start with the degree-$m$ family obtained from the preceding inverse
order. Lemma~\ref{lem:DR} makes it a degree--rank $(m,m)$ family
with the same functions. Take $A=1$, viewed as a constant unit
multidegree $(1,m)$ nilcharacter. The original positive correlation
is degree-at-most-$m-1$ correlation with the lower test equal to one.

Apply the one-step premise and the preceding recovery argument at
$R=m,m-1,\ldots,1$. There are exactly $m$ steps, a fixed standard
number. A finite sequence of positive standard density losses still
leaves a positive standard density; the same holds for positive
correlation and finite complexity bounds. Products of the successive
multidegree $(1,m)$ nilcharacters retain that bound on their fixed
product model.

At the last step the residual family has support at most $(m,0)$.
Its extra degree--rank axiom says $G_{(m,0)}=G_{(m,1)}$.
But $(m,1)>(m,0)$, so the latter group is trivial. All groups of
larger degree are also trivial. Thus the ordinary filtration
$(G_{(i,0)})_i$ has degree at most $m-1$. The residual nilcharacter,
viewed merely as a bounded nilsequence, can now be multiplied into
the already present ordinary degree-at-most-$m-1$ correlation test.
The product degree is still at most $m-1$.

We obtain one internal unit multidegree $(1,m)=(1,s-1)$ nilcharacter
$A_{\mathrm{final}}(h,n)$ such that $\Delta_h f$ has
degree-at-most-$s-2$ correlation with it on a dense set of shifts.
This is the exact Theorem 7.1 endpoint needed by the separately
authored symmetry proof. It is conditional on the one-step premise
being supplied at every input rank and every constructed family.
Acceptance of the finite iteration alone does not establish those
analytic premises or the full inverse theorem.

\section{Source findings and bounded related-use scope}

The actual final printed p.~1262 was inspected visually. In Example 6.11
the multidegree-to-degree group is printed with $|a|\leq i$, which
includes $G_0$ at every $i$ and does not give bounded degree. For
example $G_0=G_1=\R$, $G_2=0$ would give the group $\R$ at every
index. The required rule is $|a|\geq i$, proved in Section 2 above.

The same example's degree-to-rank rule prints the equality of summed
weights to the fixed bound $d$, rather than to the running index $d'$.
The desired identity $H_{(d',0)}=G_{d'}$ then need not hold. For
instance the linear filtration $G_0=G_1=\R$, $G_2=0$ with allowed
degree bound $d=2$ yields zero at $(1,0)$ under that printed rule,
whereas $G_1=\R$. A nonconstant linear polynomial is then lost.
Section 1 supplies the current-index rule and proves every filtration,
lattice and polynomial assertion used here. The printed length
condition is $\ell\geq\max(r',1)$; that sign is not itself an error.

The bounded related-use scope is Example 6.11, the polynomial identity
(6.8), Lemma 6.21's degree conversions, the initialization and rank-zero
termination of Theorems 7.1--7.2, and the diagonal/total-degree conversions
used in E.9 and the Section 13 assembly. Separate authors of the latter
consumers have been notified. No claim is made here to have read every
unrelated use of a filtration in the source or every historical correction.
These formula errors are not evidence that the intended inverse theorem
is false. Earlier source bytes and decisions are preserved.

\end{document}
